\section{Fricke constants as glue orders}
This section explains why the numerical constants in the Fricke involutions
are the indices of the root lattices in their Niemeier extensions.
For an even unimodular lattice \(N\) containing a full-rank root lattice
\(R\), the index formula gives
\[
 [N:R]^2=\det R.
\]
If \(R=A_{n-1}^{e}\), \(e=24/(n-1)\), then \(\det R=n^e\).
Consequently the glue order is
\[
 C_n=n^{12/(n-1)}.
\]
Put \(T_n(\tau)=(\eta(\tau)/\eta(n\tau))^e\).
The root systems and the relevant genus-zero functions are listed in
Cheng--Duncan--Harvey \cite[Sect.~2.2, Table~1]{CDH2014}.

\begin{theorem}\label{thm:fricke-glue}
For \(n\in\{2,3,4,5,7,9,13,25\}\), the normalized Hauptmodul of
the group generated by \(\Gamma_0(n)\) and the Fricke involution \(W_n\) is
\[
 J_n^+=T_n+e+\frac{C_n}{T_n}.
\]
The corresponding Monster classes, in the same order, are
\(2A,3A,4A,5A,7A,9A,13A,25A\).
For the root system \(A_5^4D_4\), put
\[
 T_6=\frac{\eta(\tau)^5\eta(3\tau)}
              {\eta(2\tau)\eta(6\tau)^5}.
\]
The glue order is \(72\), and the normalized Hauptmodul
\[
 J_6=T_6+5+\frac{72}{T_6}
\]
is the McKay--Thompson series \(6B\). Its group is
\(\langle\Gamma_0(6),W_6\rangle\).
\end{theorem}
\begin{proof}
The transformation \(\eta(-1/\tau)=\sqrt{-i\tau}\eta(\tau)\)
gives \(T_n(W_n\tau)=C_n/T_n(\tau)\).
The eta congruences and square character show that \(T_n\) is a
modular function on \(X_0(n)\). At a cusp with denominator \(c\mid n\)
its order is
\[
 \frac{e}{24}\frac{n}{\gcd(c^2,n)}
                   \left(1-\frac{c^2}{n}\right).
\]
For every \(n\) in the statement this is \(1\) at \(c=1\),
\(-1\) at \(c=n\), and zero at any intermediate denominator.
There are no zeros or poles in the upper half-plane, since eta
does not vanish there. Thus the divisor is \((0)-(\infty)\);
a function with a unique simple pole has degree one and identifies
\(X_0(n)\) with the projective line. Its expansion begins
\(q^{-1}-e+O(q)\).
For an indeterminate \(t\), the fixed field of
\(t\mapsto C_n/t\) in \(\C(t)\) is \(\C(t+C_n/t)\):
\(t\) satisfies \(t^2-(t+C_n/t)t+C_n=0\), and the
involution is nontrivial. Thus the displayed invariant generates
the function field of the quotient and has expansion \(q^{-1}+O(q)\).
This proves the normalized Hauptmodul assertion.

For \(A_5^4D_4\), the determinant is \(6^4\cdot4\), with square root
\(72\). Applying the same eta transformation to the four factors
of \(T_6\) gives \(T_6(W_6\tau)=72/T_6(\tau)\).
The cusp orders of \(T_6\), at denominators \(1,2,3,6\), are
\(1,0,0,-1\). The same degree-one argument makes \(T_6\) a
coordinate on \(X_0(6)\).
Here \(T_6=q^{-1}-5+6q+4q^2+O(q^3)\), so the same fixed-field
argument proves the formula for \(J_6\).
The Monster identifications, including the Fricke-only group \(6+6\)
for \(6B\), are those of Conway--Norton
\cite[Table~2, pp.~327--328]{CN1979}.
At level \(25\) one must distinguish the Fricke quotient \(25A\)
from the function for \(\Gamma_0(25)\), traditionally called \(25Z\).
\end{proof}

\begin{center}
\begin{tabular}{c|rrrrrrrr}
\(n\)&2&3&4&5&7&9&13&25\\ \hline
\(e\)&24&12&8&6&4&3&2&1\\
\(C_n\)&4096&729&256&125&49&27&13&5
\end{tabular}
\end{center}

\section{The cusp cover of \(X_0(6)\)}
The four branch values of the cusp cover give a genus-one curve
whose relation with the conductor-\(24\) modular curve can be proved
by rational functions.

\begin{theorem}\label{thm:cusp-isogeny}
The smooth projective completion of
\[
 C:\quad y^2=t(t+8)(t+9),
\]
with its point at infinity as origin, admits a degree-two isogeny to
\[
 E:\quad v^2=(x-1)(x-2)(x+2)=x^3-x^2-4x+4.
\]
On the affine open set \(t\ne-9\), the isogeny is
\begin{equation}\label{eq:cusp-isogeny}
 x=\frac{t^2+16t+72}{4(t+9)},\qquad
 v=\frac{y(t+6)(t+12)}{8(t+9)^2}.
\end{equation}
Its kernel is \(\{O,(-9,0)\}\), and its pullback of \(dx/(2v)\)
is \(2\,dt/(2y)\). The target is Cremona's curve \(24a1\),
whose LMFDB label is \(24.a4\). Thus the Hasse--Weil \(L\)-function
of \(C\) is that of \(\eta(2\tau)\eta(4\tau)\eta(6\tau)\eta(12\tau)\),
including all finite local factors.
\end{theorem}
\begin{proof}
Set \(u=t+9\). Then \(C\) becomes
\[
 y^2=u^3-10u^2+9u.
\]
For a curve \(y^2=u^3+au^2+bu\), direct substitution shows that
\[
 X=u+a+\frac bu,\qquad Y=y\left(1-\frac b{u^2}\right)
\]
satisfy
\[
 Y^2=X^3-2aX^2+(a^2-4b)X.
\]
With \(a=-10,b=9\), this is \(Y^2=X^3+20X^2+64X\).
The change \(X=4(x-2),Y=8v\) gives \(E\) and
the functions in \eqref{eq:cusp-isogeny}.
The map extends over the smooth projective curves. It sends the origin
to the origin, so is a group homomorphism. The rational function \(X\)
has degree two as a function of \(u\); equivalently,
\(u^2+(a-X)u+b=0\), and the kernel is precisely the indicated
two-torsion subgroup. This proves the degree and kernel.
Differentiating the two formulas gives the differential assertion.

For completeness, the dual map, in the same coordinates, is
\[
 t=x-6+\frac4{x-2},\qquad
 y=v\left(1-\frac4{(x-2)^2}\right).
\]
The two compositions are multiplication by two.
The equation and the conductor-\(24\) identification of \(E\)
are recorded in \cite{LMFDB24}. The eta quotient in the statement
has weight two, trivial character and positive orders at every
cusp of \(\Gamma_0(24)\); the two eta congruence sums equal \(24\),
and the product entering the character is \(576\).
The space \(S_2(\Gamma_0(24))\) has dimension one and its old space
is zero. The quotient is therefore the normalized newform of \(E\).
An isogeny identifies the rational Tate modules, and hence their
inertia-invariant Frobenius characteristic polynomials at every prime.
This gives equality of all finite local factors.
\end{proof}

\section{CM zeros of Niemeier theta series}
The zeros at imaginary quadratic arguments are decided by the rational
singular moduli, rather than by a Riemann hypothesis for a Mellin transform.
Write
\(\Theta_L(\tau)=\sum_{v\in L}q^{(v,v)/2}\), let \(N_2(L)\) be the
number of roots, and set \(h=N_2(L)/24\); for the Leech lattice
\(\Lambda\), this convention gives \(h=0\).

\begin{theorem}\label{thm:cm-theta-zeros}
Among the \(24\) positive-definite even unimodular lattices of rank \(24\),
\(\Theta_L\) vanishes at an imaginary quadratic argument only for
the root systems \(E_8^3\) and \(D_{16}E_8\). In both cases the
zero arguments are exactly the \(\mathrm{SL}_2(\Z)\)-orbit of
\(\rho=e^{2\pi i/3}\).
The genus-two theta series of the Leech lattice does not vanish at
any point on the Humbert surface \(H_1\) represented by a pair
of imaginary quadratic arguments.
\end{theorem}
\begin{proof}
The two-dimensional space of weight-twelve modular forms on
\(\mathrm{SL}_2(\Z)\) gives
\begin{equation}\label{eq:niemeier-j}
 \Theta_L=E_4^3+(24h-720)\Delta
          =\Delta\bigl(j-720+24h\bigr),
\end{equation}
by comparison of the constant and \(q\) coefficients.
The product formula for \(\Delta\) shows that it has no zero in
the upper half-plane. The possible Coxeter numbers are
\[
 0,2,3,4,5,6,7,8,9,10,12,13,14,16,18,22,25,30,46.
\]
Their target values \(720-24h\) lie in \([-384,720]\).
If \(j(\tau)\) is rational at an imaginary quadratic argument,
the class polynomial has degree one. The complete list of these
rational singular moduli is given in
\cite[Table~1, p.~1057]{BLM2017}; the only member in this interval
is \(0\), of discriminant \(-3\).
Hence \(h=30\), and the Niemeier classification gives precisely
the two stated root systems. The equation \(j(\tau)=0\)
characterizes the modular orbit of \(\rho\).

On a block diagonal period matrix,
the definition of the genus-two theta series gives
\[
 \vartheta_2(L)(\operatorname{diag}(\tau_1,\tau_2))
       =\Theta_L(\tau_1)\Theta_L(\tau_2).
\]
For \(L=\Lambda\), neither factor is zero when both arguments
are imaginary quadratic. Every point of \(H_1\) has such a
block diagonal representative, and the Siegel modular transformation
factor is nonzero. The claimed nonvanishing follows.
This assertion concerns \(H_1\); it makes no claim away from that
surface or about the zeros of an Epstein zeta function.
\end{proof}

\section{Ogg's list and explicit class-number cutoffs}
Ogg's fixed-point formula already expresses the genus-zero question
in terms of quadratic class numbers. The additional assertion here is
an explicit finite cutoff, first elementary and then smaller under
the Dirichlet Riemann hypothesis.
Let \(h(D)\) denote the proper ideal class number of the quadratic
order of discriminant \(D<0\), and let \(\chi_D\) denote the
Kronecker character. For a prime \(p\ge5\), set
\[
 r(p)=\frac{h(-4p)+{\bf1}_{p\equiv3\ (4)}h(-p)}2,\qquad
 s(p)=\frac{p-1}{12}+\frac{1-\chi_{-4}(p)}4+
                         \frac{1-\chi_{-3}(p)}3.
\]

\begin{theorem}\label{thm:ogg-cutoffs}
For every prime \(p>7000\), one has \(r(p)<s(p)\), without any
Riemann hypothesis. Under the Dirichlet Riemann hypothesis the same
conclusion follows for every prime \(p>4300\).
The equality \(r(p)=s(p)\), for \(p\ge5\), holds exactly for
\[
 5,7,11,13,17,19,23,29,31,41,47,59,71.
\]
Together with \(2,3\), these are the primes for which
\(X_0(p)/\langle W_p\rangle\) has genus zero.
\end{theorem}
\begin{proof}
Ogg's Riemann--Hurwitz calculation
\cite[p.~452]{Ogg1974}, fixed-point formula
\cite[p.~454]{Ogg1974}, and genus formula
\cite[p.~455]{Ogg1974} give
\[
 g(X_0(p))=s(p)-1,\quad
 \#\operatorname{Fix}(W_p)=2r(p),\quad
 g(X_0(p)/\langle W_p\rangle)=\frac{s(p)-r(p)}2.
\]
Thus the class-number formulation itself is classical.

For a nonprincipal real primitive character of conductor \(q\),
periodicity and zero mean give
\[
 |L(1,\chi)|\le H_q+1\le\log q+2.
\]
Indeed truncate the Dirichlet series at \(q\), bound its initial
sum by \(H_q\), and integrate the tail by parts using
\(\left|\sum_{n\le x}\chi(n)\right|\le q\).
When \(p\equiv1\pmod4\), the class-number formula gives
\[
 r(p)=\frac{\sqrt p}{\pi}L(1,\chi_{-4p}).
\]
When \(p\equiv3\pmod4\), the order-conductor formula gives
\[
 h(-4p)=(2-\chi_{-p}(2))h(-p),\qquad
 r(p)\le \frac{2\sqrt p}{\pi}L(1,\chi_{-p}).
\]
Since \(\log(4p)+2\le2(\log p+2)\), both cases imply
\begin{equation}\label{eq:ogg-elementary}
 r(p)\le\frac{2\sqrt p}{\pi}(\log p+2).
\end{equation}
At \(p=7000\),
\[
 \frac{p-1}{\sqrt p}-\frac{24}{\pi}(\log p+2)>0.73815.
\]
The derivative of the left side is positive for \(p\ge7000\):
it is at least \(1/(2\sqrt p)-24/(\pi p)>0\).
Consequently \eqref{eq:ogg-elementary} is strictly smaller
than \((p-1)/12\le s(p)\) beyond that cutoff.

For the smaller cutoff, put
\[
 U(q)=2e^\gamma
 \left(\log\log q-\log2+\frac12+\frac1{\log\log q}\right).
\]
Under the Dirichlet Riemann hypothesis, the explicit estimate of
Lamzouri--Li--Soundararajan, extended to every \(q\ge3\) by
Languasco--Trudgian \cite[Theorem~1]{LT2022}, gives
\(L(1,\chi)\le U(q)\).
For \(p\ge4300\), \(U(4p)\le2U(p)\).
To see this, put \(x=\log\log p>1\).
The difference \(\log\log(4p)-x\) is at most \(\log2\);
the function \(x-\log2+1/2+1/x\) has derivative between zero
and one there, and its value is greater than \(\log2\).
It follows that \(r(p)\le2\sqrt p\,U(p)/\pi\).
At \(p=4300\),
\[
 \frac{p-1}{\sqrt p}-\frac{24}{\pi}U(p)>0.19855.
\]
Its derivative is at least
\[
 \frac1{2\sqrt p}-\frac{48e^\gamma}{\pi p\log p}>0
 \qquad(p\ge4300).
\]
This proves the second cutoff. Both numerical inequalities can
equivalently be checked with rational interval bounds for
\(\pi,e^\gamma,\log7000,\log4300\).

Finally, enumerate primitive reduced positive binary quadratic forms
\([a,b,c]\), with
\[
 b^2-4ac=D,\quad |b|\le a\le c,\quad
 b\ge0\ \text{if } |b|=a\text{ or }a=c.
\]
The reduction inequalities give \(a\le\sqrt{|D|/3}\), so this is
a finite exact computation. For each of the \(898\) primes
\(5\le p\le7000\), counting these forms for the two discriminants
in \(r(p)\) gives the displayed equality list and no instance of
\(r(p)>s(p)\). The workbench program \emph{Root exact checks}
records every integer count; a separate enumeration gives the
same list. This completes the finite verification.
\end{proof}
